\section{Functions: Structured Promises, Not Just Formulas}
\chapterline

In elementary algebra and calculus, a function is almost always introduced as an algebraic formula: an expression like $f(x) = x^2$ or $g(x) = \sin(x) / x$. You plug a real number into the machine, turn the symbolic crank, and read off the computed result.

In pure mathematics, that formula-centric perspective is entirely too narrow. A formula is merely an algorithm or assignment rule. A \emph{function} is a structured contract between two mathematical sets. It is a promise that transforms an entire universe of inputs into a declared destination space, obeying strict geometric and logical rules.

\begin{historybox}
The modern definition of a function was forged through a century of fierce mathematical debate. In 1718, Johann Bernoulli defined a function as ``any quantity composed in whatever manner of a variable and constants.'' Thirty years later, Leonhard Euler (\emph{Introductio in analysin infinitorum}, 1748) codified functions as expressions formed by algebraic operations or infinite series. 

This formula-based view triggered the famous 18th-century ``vibrating string controversy'' between Jean le Rond d'Alembert, Euler, and Daniel Bernoulli. If a plucked violin string is released from an initial triangular shape with a sharp corner, can that physical motion be described by a function? Euler argued that a curve with a corner could not be a genuine function because it had no single continuous differential equation. 

The breakthrough came in 1829 from Peter Gustav Lejeune Dirichlet while studying the convergence of Fourier series. Dirichlet realized that mathematics had to sever the concept of a function from the demand for an algebraic formula. A function, Dirichlet proposed, is simply an arbitrary assignment that pairs every input $x$ with a unique output $y$. To prove how far this extended beyond formulas, he introduced the notorious Dirichlet function:
\[
\mathbf{1}_\mathbb{Q}(x) = \begin{cases} 1 & \text{if } x \in \mathbb{Q}, \\ 0 & \text{if } x \notin \mathbb{Q}. \end{cases}
\]
Dirichlet's function has no algebraic formula, cannot be drawn with a pen, and is discontinuous at every single real number! Traditionalists were horrified; Charles Hermite famously lamented, ``I turn away with fright and horror from this lamentable plague of functions which have no derivatives.'' Yet Dirichlet was right. Richard Dedekind and later Nicolas Bourbaki (1939) cemented this insight in modern set theory: a function is not an algebraic equation, but a lawful subset of a Cartesian product.
\end{historybox}

\subsection{The Anatomy of a Function}

Formally, a function consists of three distinct components:
\begin{enumerate}[label=\arabic*.]
\item A set $A$ called the \textbf{domain} (the universe of allowed inputs).
\item A set $B$ called the \textbf{codomain} (the designated target space).
\item A relation or rule that assigns to \textbf{every} element $a \in A$ \textbf{exactly one} element $b \in B$, denoted $f(a) = b$.
\end{enumerate}

We package this entire tripartite package in the standard notation:
\[
f: A \to B.
\]

\begin{definition}[Function]
A function $f: A \to B$ is a relation from $A$ to $B$ satisfying two structural conditions:
\begin{enumerate}[label=(\roman*)]
\item \textbf{Total Definition:} For every $a \in A$, there exists some $b \in B$ such that $f(a) = b$.
\item \textbf{Single-Valuedness:} For every $a \in A$ and all $b_1, b_2 \in B$, if $f(a) = b_1$ and $f(a) = b_2$, then $b_1 = b_2$.
\end{enumerate}
\end{definition}

\begin{center}
\begin{tikzpicture}[scale=0.92, >=Latex]
  % Lane 1: A Valid Function
  \begin{scope}[yshift=0cm]
    \node[anchor=west, font=\small\bfseries] at (-0.5, 3.2) {A Valid Function $f: X \to Y$};
    \draw[rounded corners=4pt, fill=gray!3, draw=gray!25] (-0.6, -0.6) rectangle (12.2, 3.0);
    
    % Domain X
    \draw[thick, sourceblue, fill=sourceblue!10] (1.3, 1.2) ellipse (1.1cm and 1.4cm);
    \node[sourceblue!90!black, font=\footnotesize\bfseries] at (1.3, 2.75) {Domain $X$};
    
    % Codomain Y
    \draw[thick, targetgreen, fill=targetgreen!8] (4.8, 1.2) ellipse (1.4cm and 1.4cm);
    \node[targetgreen!90!black, font=\footnotesize\bfseries] at (4.8, 2.75) {Codomain $Y$};
    
    % Image subset
    \draw[thick, dashed, quotientorange, fill=quotientorange!15] (4.6, 1.05) ellipse (0.95cm and 0.95cm);
    \node[quotientorange!90!black, font=\scriptsize\bfseries] at (4.6, -0.15) {Image $\operatorname{im}(f)$};
    
    % Points in X
    \node[circle, fill=sourceblue, inner sep=1.8pt, label={left:\footnotesize $x_1$}] (x1) at (1.3, 1.9) {};
    \node[circle, fill=sourceblue, inner sep=1.8pt, label={left:\footnotesize $x_2$}] (x2) at (1.3, 1.2) {};
    \node[circle, fill=sourceblue, inner sep=1.8pt, label={left:\footnotesize $x_3$}] (x3) at (1.3, 0.5) {};
    
    % Points in Y
    \node[circle, fill=quotientorange!80!black, inner sep=1.8pt, label={right:\footnotesize $y_1$}] (y1) at (4.4, 1.55) {};
    \node[circle, fill=quotientorange!80!black, inner sep=1.8pt, label={right:\footnotesize $y_2$}] (y2) at (4.7, 0.65) {};
    \node[circle, fill=gray!60!black, inner sep=1.8pt, label={right:\footnotesize $y_3$}] (y3) at (5.5, 1.75) {};
    
    % Arrows
    \draw[->, thick, sourceblue!85!black] (x1) -- (y1);
    \draw[->, thick, sourceblue!85!black] (x2) -- (y1);
    \draw[->, thick, sourceblue!85!black] (x3) -- (y2);
    
    \node[anchor=west, font=\scriptsize, align=left] at (6.6, 1.2) {
      $\bullet$ Every input emits \textbf{exactly one} beam.\\
      $\bullet$ Arrows may converge ($x_1, x_2 \mapsto y_1$).\\
      $\bullet$ Codomain elements ($y_3$) may be unhit.
    };
  \end{scope}

  % Lane 2: Broken Contracts (Non-Functions)
  \begin{scope}[yshift=-4.6cm]
    \node[anchor=west, font=\small\bfseries] at (-0.5, 3.2) {Two Ways to Break the Contract (Not Functions)};
    \draw[rounded corners=4pt, fill=trapred!3, draw=trapred!25] (-0.6, -0.6) rectangle (12.2, 3.0);
    
    % Panel A: Missing an output
    \begin{scope}[xshift=0cm]
      \draw[thick, trapred!70, fill=white] (1.2, 1.2) ellipse (0.9cm and 1.3cm);
      \node[trapred!90!black, font=\scriptsize\bfseries] at (1.2, 2.2) {Domain $A$};
      \draw[thick, gray!50, fill=white] (4.0, 1.2) ellipse (0.9cm and 1.3cm);
      \node[gray!70!black, font=\scriptsize\bfseries] at (4.0, 2.2) {Codomain $B$};
      
      \node[circle, fill=sourceblue, inner sep=1.6pt, label={left:\scriptsize $u$}] (u) at (1.2, 1.7) {};
      \node[circle, fill=trapred, inner sep=2pt, label={left:\scriptsize\bfseries $v$}] (v) at (1.2, 0.7) {};
      \node[circle, fill=gray!60, inner sep=1.6pt] (w) at (4.0, 1.5) {};
      
      \draw[->, thick, sourceblue!80] (u) -- (w);
      \node[below, trapred!80!black, font=\scriptsize\bfseries] at (2.6, -0.2) {Missing: $v$ has no output};
    \end{scope}
    
    % Panel B: Ambiguous multiple outputs
    \begin{scope}[xshift=6.5cm]
      \draw[thick, trapred!70, fill=white] (1.2, 1.2) ellipse (0.9cm and 1.3cm);
      \node[trapred!90!black, font=\scriptsize\bfseries] at (1.2, 2.2) {Domain $A$};
      \draw[thick, gray!50, fill=white] (4.0, 1.2) ellipse (0.9cm and 1.3cm);
      \node[gray!70!black, font=\scriptsize\bfseries] at (4.0, 2.2) {Codomain $B$};
      
      \node[circle, fill=trapred, inner sep=2pt, label={left:\scriptsize\bfseries $p$}] (p) at (1.2, 1.2) {};
      \node[circle, fill=gray!60, inner sep=1.6pt, label={right:\scriptsize $q_1$}] (q1) at (4.0, 1.8) {};
      \node[circle, fill=gray!60, inner sep=1.6pt, label={right:\scriptsize $q_2$}] (q2) at (4.0, 0.6) {};
      
      \draw[->, thick, trapred!80] (p) -- (q1);
      \draw[->, thick, trapred!80] (p) -- (q2);
      \node[below, trapred!80!black, font=\scriptsize\bfseries] at (2.6, -0.2) {Ambiguous: $p$ splits into two};
    \end{scope}
  \end{scope}
\end{tikzpicture}
\end{center}

\begin{intuitionbox}
Think of the domain as an array of laser pointers, and the codomain as a screen across the room. A valid function requires every single laser pointer to be switched on, and each pointer must project a clean, single dot on the screen. It is completely acceptable for two different lasers to shine on the same spot, or for parts of the screen to remain unlit. But a laser cannot be left dark, and a laser cannot shoot a beam that splits into two separate targets.
\end{intuitionbox}

\subsection{Why the Formula Alone Is Incomplete}

Consider the squaring formula $x \mapsto x^2$. Stating only the formula leaves out the actual behavior of the function. Watch how drastically the object transforms as we alter the domain and codomain:

\begin{itemize}
\item \textbf{Function 1:} $f: \R \to \R$ with $f(x) = x^2$. 
Because negative numbers produce positive squares, $f(-2) = f(2) = 4$. Multiple inputs collide on the same output. Furthermore, negative numbers in the codomain like $-7$ never get hit.
\item \textbf{Function 2:} $g: \R \to [0, \infty)$ with $g(x) = x^2$.
Here, the codomain has been trimmed to match the actual outputs. Every target $y \ge 0$ is reached by at least one input ($x = \sqrt{y}$), but collisions still occur ($g(-2) = g(2)$).
\item \textbf{Function 3:} $h: [0, \infty) \to [0, \infty)$ with $h(x) = x^2$.
Now negative inputs are eliminated from the domain. Collisions are gone, and every target in the codomain is hit. This function has an exact two-way inverse: $h^{-1}(y) = \sqrt{y}$.
\end{itemize}

Even though all three functions share the identical formula $x^2$, $f$ is neither invertible nor surjective, $g$ is surjective but not injective, and $h$ is a perfect bijection. Domain and codomain are not incidental metadata; they define what the function is.

\subsection{Codomain Versus Image}

One of the most persistent hurdles for beginners is confusing the \emph{codomain} with the \emph{image} (often called the range).

\begin{definition}[Image]
Let $f: A \to B$ be a function. The \textbf{image} of $f$, denoted $\operatorname{im}(f)$ or $f(A)$, is the subset of the codomain consisting of values that are actually achieved:
\[
\operatorname{im}(f) = \{f(a) \in B : a \in A\} = \{b \in B : (\exists a \in A)(f(a) = b)\}.
\]
\end{definition}

The codomain $B$ is the room you chose to throw darts in; the image $\operatorname{im}(f)$ is the collection of dart holes on the wall. The image is always a subset of the codomain:
\[
\operatorname{im}(f) \subseteq B.
\]
When every point in the codomain receives a dart ($\operatorname{im}(f) = B$), the function is called \emph{surjective}.

\subsection{Connection to Software Engineering}

In computer science, a function's domain and codomain correspond directly to its type signature:
\begin{center}
\texttt{function processData(x: InputType): ReturnType}
\end{center}
Declaring \texttt{ReturnType} as \texttt{int64} does not mean the function produces every possible 64-bit integer; it guarantees that whatever it returns will live within that set. A bug where a function throws an unhandled exception for some input is a failure of total definition. A function that behaves non-deterministically without explicit random state violates single-valuedness.

\subsection{What Modern Functions Unlocked, Their Limits, and Where We Go Next}

Freeing functions from formulas unlocked vast domains of modern science. It gave mathematicians the power to define discrete state-machine transitions, characteristic indicator functions in probability, piecewise step voltages in signal processing, and cryptographic permutations—none of which can be written as smooth algebraic expressions. It transformed mathematics from a catalog of formulas into a science of lawful transformations.

Yet knowing that an assignment $f: A \to B$ is a valid function tells you surprisingly little about the quality of the transformation:
\begin{itemize}
\item \textbf{Information loss:} A function is entirely permitted to map every single element of a million-element domain to the number $0$. It satisfies the definition of a function perfectly, but it completely obliterates the structure of the input space.
\item \textbf{Coverage gaps:} The codomain could be the entire real line $\R$, while the actual outputs form only a tiny handful of numbers.
\item \textbf{Irreversibility:} If multiple inputs collide on the same output, you can never trace a result backward to find the original cause.
\end{itemize}

A valid function is only a raw, minimal promise. To know whether a mapping preserves information, covers its destination, or can be undone, we need sharper structural tools. In Chapter~13, we introduce \textbf{Injective, Surjective, Bijective}, establishing the three fundamental pillars that govern collision-free mappings, target coverage, and invertible pairs.

\subsection{Exercises}

\begin{enumerate}
\item Determine whether each proposed rule defines a valid function. If not, state which property fails:
\begin{enumerate}[label=(\alph*)]
\item $f: \R \to \R$ defined by $f(x) = \frac{1}{x^2 - 4}$.
\item $g: \Q \to \Z$ defined by $g(p/q) = p + q$.
\item $h: \R \to \R$ defined by $h(x) = y$ where $y^2 = x$.
\end{enumerate}
\item Find the image of the function $f: \R \setminus \{1\} \to \R$ given by $f(x) = \frac{2x+1}{x-1}$. Is the number $2$ in the image?
\item Let $A = \{1, 2, 3\}$ and $B = \{a, b\}$. How many distinct functions exist from $A$ to $B$? How many distinct functions exist from $B$ to $A$?
\end{enumerate}

\newpage
